% =====================================================================
% المقدمة العامة وإشكالية الدراسة
% =====================================================================

\chapter*{مقدمة عامة وإشكالية الدراسة}
\addcontentsline{toc}{chapter}{مقدمة عامة وإشكالية الدراسة}
\markboth{مقدمة عامة وإشكالية الدراسة}{مقدمة عامة وإشكالية الدراسة}

%% ---------------------------------------------------------------------
\section*{تمهيد وتأصيل تاريخي واقتصادي}
%% ---------------------------------------------------------------------

شهدت المنظومة الاقتصادية والإنتاجية العالمية خلال العقود الأخيرة تحولات جذرية فرضتها عولمة الأسواق، واحتدام المنافسة، وثورة التكنولوجيا وسرعة تدفق رؤوس الأموال. وقد ألقت هذه التحولات بظلالها المباشرة على العلاقات المهنية وأنماط إدارة الموارد البشرية، حيث تراجعت المقاولات الحديثة تدريجياً عن النمط التقليدي للتشغيل القائم على الاستقرار الوظيفي وعقد الشغل الفردي غير المحدد المدة (CDI) كقاعدة عامة وحيدة، لصالح البحث المحموم عما بات يعرف في الأدبيات الاقتصادية والقانونية بـ \textbf{«مرونة الشغل»} (\textit{Flexibilité de l'emploi}).

وتتجلى هذه المرونة في قدرة المنشأة على ملاءمة حجم وقدرات مواردها البشرية مع تقلبات السوق ودورات الإنتاج وفترات الذروة المفاجئة دون تحمل الأعباء المالية والتنظيمية الثقيلة المترتبة على إنهاء عقود الشغل القارة. وفي هذا السياق، برز نمطان أساسيان ومترابطان لإعادة هيكلة قوة العمل:
\begin{enumerate}
    \item \textbf{الشغل المؤقت (\textit{Travail temporaire / Intérim}):} وهو علاقة ثلاثية الأبعاد تقوم بمقتضاها مقاولة متخصصة مرخص لها بوضع أجراء ترتبط معهم بروابط تعاقدية رهن إشارة منشأة أخرى مستعملة، لأداء مهام غير دائمة وبشروط محددة حصراً.
    \item \textbf{المناولة أو المقاولة من الباطن (\textit{Sous-traitance}):} وهي تقنية تعاقدية تجارية يعهد بموجبها مقاول أصلي إلى مقاول آخر مستقل بإنجاز عمل أو تقديم خدمة متخصصة، محتفظاً بالتبعية المهنية والفنية لأجرائه.
\end{enumerate}

وعلى صعيد النسيج الاقتصادي المغربي، تسارع اللجوء إلى هاتين الآليتين منذ أواخر تسعينيات القرن الماضي مع انفتاح الاقتصاد الوطني وتوقيع اتفاقيات التبادل الحر وتنامي قطاعات استراتيجية كالسيارات والطيران واللوجستيك والصناعات الغذائية. ورغم أن هذا التوجه وفّر متنفساً حقيقياً للمقاولات لتخفيض تكاليف التسيير وضمان التنافسية، وفتح أمام فئات واسعة من الشباب والباحثين عن شغل بوابة للاندماج المهني واكتساب المهارات، إلا أنه أفرز في المقابل إشكالات بنيوية بالغة الخطورة؛ حيث أدى تفكيك الرابطة العقدية الثنائية المباشرة بين الأجير ورب العمل الفعلي إلى إضعاف الحماية الاجتماعية، وتشتيت عناصر التبعية القانونية، وتنامي مظاهر الهشاشة المهنية وتدني الأجور.

%% ---------------------------------------------------------------------
\section*{الإطار الدستوري والمرجعية الدولية}
%% ---------------------------------------------------------------------

تتقاطع إشكالية الشغل المؤقت والمناولة مع ثنائية دستورية بالغة الدقة؛ فمن جهة، يكرس \textbf{دستور المملكة المغربية لسنة 2011} في فصله الخامس والثلاثين (35) حرية المبادرة والمقاولة، والتنافس الحر، وضمان مناخ ملائم للاستثمار. ومن جهة ثانية، ينص الفصل الحادي والثلاثون (31) على التزام الدولة والمؤسسات العمومية والجماعات الترابية بتعبئة كل الوسائل المتاحة لتيسير استفادة المواطنات والمواطنين من الحق في الشغل، وضمان شروط عمل لائقة وعادلة، والاستفادة من الحماية الاجتماعية والتغطية الصحية والتكوين المهني.

ودولياً، تكتسي هذه المسألة عناية خاصة لدى \textbf{منظمة العمل الدولية (OIT)}، التي وضعت معايير متقدمة للتوفيق بين مرونة السوق وحماية العمال:
\begin{itemize}
    \item \textbf{الاتفاقية رقم 181 بشأن وكالات الاستخدام الخاصة (1997):} التي كرست الاعتراف بالدور الإيجابي لشركات التشغيل المؤقت في تسيير سوق الشغل، شريطة إخضاعها لترخيص صارم، ومنع اقتطاع أي مصاريف أو أتعاب من أجور العمال، وإقرار مسؤولية مشتركة واضحة تضمن تمتع الأجراء بكافة حقوقهم الأساسية والتفاوض الجماعي.
    \item \textbf{التوصية رقم 198 بشأن علاقة العمل (2006):} التي شددت على المعيار الجوهري المتمثل في \textbf{«مبدأ أولوية الواقع على المظاهر والشكل»} (\textit{Principe de la primauté des faits})؛ حيث تقضي بوجوب تكييف علاقة الشغل بالنظر إلى الوقائع الميدانية وممارسة السلطة الفعلية، بصرف النظر عن التسميات أو العقود الصورية التي يضفيها الأطراف للتنصل من موجبات القانون.
\end{itemize}

%% ---------------------------------------------------------------------
\section*{سياق الدراسة وأهميتها}
%% ---------------------------------------------------------------------

تنطلق هذه الدراسة من واقعة عملية ذات راهنية قصوى تتكرر بصورة شبه نمطية في العديد من المنشآت الصناعية والإنتاجية الكبرى بالمغرب:
\begin{quote}
\textit{«شركة صناعية كبرى تعقد اتفاقاً مع مقاولة متخصصة تحت غطاء "المناولة"، غير أن المحتوى الفعلي للعقد ينحصر في توفير اليد العاملة والعمالة اليدوية. ويشتغل هؤلاء العمال داخل أسوار المصنع ذاته لسنوات متتالية، جنباً إلى جنب مع أجراء المقاولة الأصلية، مؤدين نفس المهام والوظائف المندرجة في النشاط العادي والمستمر للمصنع، وتحت الإشراف والتوجيه المباشر لأطر المقاولة المستفيدة. ورغم ذلك، يتقاضى هؤلاء العمال أجوراً أدنى، ويحرمون من التعويضات والترقيات والامتيازات الاجتماعية الممنوحة لنظرائهم القارين، وسط غموض كامل يكتنف الجهة المسؤولة قانونياً عن سلامتهم وحمايتهم وضمان حقوقهم».}
\end{quote}

وتتفرع أهمية تفكيك هذه النازلة إلى أربعة أبعاد متداخلة:
\begin{itemize}
    \item \textbf{البعد القانوني والتنظيمي:} ضبط الفوارق الماهوية الفاصلة بين الشغل المؤقت، والمقاولة من الباطن، وعقد الشغل المحدد المدة (CDD)، والتمييز بين الإسناد المشروع للخدمات وبين جريمة "إعارة اليد العاملة غير المشروعة بمقابل" والتحايل على أحكام مدونة الشغل (القانون رقم 65.99).
    \item \textbf{البعد الاجتماعي والإنساني:} صيانة كرامة الأجراء ودرء الهشاشة المزدوجة الناتجة عن تدني الأجر وعدم استقرار الرابطة العقدية، وتكريس المبدأ الحقوقي الكوني: «أجر متساوٍ عن عمل متساوي القيمة»، وضمان المساواة التامة في شروط الوقاية والسلامة وحوادث الشغل.
    \item \textbf{البعد الاقتصادي والتنافسي:} تحديد التخوم الفاصلة بين المرونة الاقتصادية المشروعة التي تعزز تنافسية المقاولة وتساعدها على امتصاص الأزمات الظرفية، وبين "الإغراق الاجتماعي" (\textit{Dumping social}) القائم على ضغط التكاليف عبر استغلال الثغرات التشريعية، بما يضر بالتنافسية الشريفة بين المقاولات.
    \item \textbf{البعد العملي والمؤسساتي:} رسم معالم المسؤولية القانونية المدنية والجنائية للأطراف (المشغل الأصلي والوسيط)، وتحديد نطاق التدخل الرقابي المتاح لمفتشية الشغل والقضاء الاجتماعي لإنفاذ مقتضيات النظام العام الاجتماعي وحماية الطرف الضعيف في العلاقة.
\end{itemize}

%% ---------------------------------------------------------------------
\section*{الإشكالية المركزية}
%% ---------------------------------------------------------------------

في ضوء التجاذب الدائم بين ضرورات المرونة التدبيرية للمقاولات من جهة، ومقتضيات الاستقرار المهني والأمن التعاقدي للأجراء من جهة ثانية، تتمحور الإشكالية المركزية لهذه الدراسة حول التساؤل الجوهري الآتي:

\begin{tcolorbox}[colback=lightgray, colframe=maincolor, arc=5pt, title=\textbf{الإشكالية المحورية}, fonttitle=\bfseries]
\centering
\large
إلى أي حد يوفق الإطار القانوني المغربي في تحقيق التوازن الحرج بين تلبية متطلبات المقاولات في المرونة وتقليص التكاليف عبر الشغل المؤقت والمناولة، وبين صيانة الحقوق الجوهرية للأجراء وحمايتهم من مظاهر الاستغلال والتحايل التعاقدي على أرض الواقع؟
\end{tcolorbox}

وتتفرع عن هذه الإشكالية الأم جملة من التساؤلات الفرعية النسقية:
\begin{enumerate}
    \item ما هي الخصائص المفاهيمية المميزة لكل من الشغل المؤقت وعقد المقاولة من الباطن، وما هي المعايير الضابطة للفصل بينهما وبين عقود الشغل محددة المدة (CDD)؟
    \item ما مدى دقة الحالات الحصرية والقيود الزمنية والموضوعية التي رسمها المشرع المغربي في المواد 495 إلى 506 من مدونة الشغل لتأطير الاستعانة بمقاولات التشغيل المؤقت؟
    \item ما هي الالتزامات والضمانات المقررة لصالح الأجير تجاه كل من المشغل القانوني (الوسيط) والمستعمل الفعلي، لا سيما في مجالات الأجر، والتصريح بالضمان الاجتماعي، والتأمين ضد حوادث الشغل والأمراض المهنية؟
    \item ما هي حدود المسؤولية القانونية الملقاة على المقاول الأصلي في عقود المناولة وفقاً للمادتين 89 و90 من المدونة، وما مدى كفاية شروط هذه المسؤولية لضمان حقوق العمال في حالة إعسار المقاول من الباطن؟
    \item هل يملك القضاء الاجتماعي وأجهزة تفتيش الشغل الأدوات والآليات الكفيلة بزجر الممارسات الصورية، وإعادة تكييف عقود المناولة التي تستعمل كغطاء لتوفير يد عاملة دائمة إلى عقود شغل مباشرة وغير محددة المدة؟
\end{enumerate}

%% ---------------------------------------------------------------------
\section*{فرضيات البحث}
%% ---------------------------------------------------------------------

سعياً للإحاطة بمختلف جوانب الموضوع واختبار صحة الطروحات الفقهية والقضائية، تنطلق دراستنا من الفرضيات العلمية الأربع التالية:
\begin{enumerate}
    \item \textbf{الفرضية الأولى (التباين الماهوي في محل الالتزام):} يختلف عقد الشغل المؤقت جوهرياً عن عقد المناولة من حيث موضوع الالتزام؛ فالشغل المؤقت ينصب حصراً على وضع طاقة عمل بشرية رهن إشارة المستعمل مع انتقال جزء من التبعية القانونية، في حين تنصب المناولة على التزام المقاول من الباطن بتحقيق نتيجة مستقلة (شغل أو خدمة) عبر وسائله وإشرافه التقني المباشر دون تفويت التبعية.
    \item \textbf{الفرضية الثانية (عدم مشروعية تغطية الحاجيات الدائمة):} إن استخدام الأجراء المؤقتين أو عمال المناولة لشغل مناصب قارة ولأداء مهام تدخل في صميم النشاط الدائم والعادي للمقاولة المستعملة لسنوات متواصلة، يمثل خرقاً جوهرياً للنظام العام الاجتماعي، ويبرر قانوناً وقضائياً إسقاط الحجاب التعاقدي وإعادة تكييف العلاقة إلى عقد شغل غير محدد المدة (CDI) مباشر مع المقاولة المستعملة.
    \item \textbf{الفرضية الثالثة (قصور نظام المسؤولية الاحتياطية بالمناولة):} يعاني نظام المسؤولية التضامنية أو الاحتياطية المنصوص عليه في المادة 89 من مدونة الشغل من ثغرات هيكلية حادة؛ لكونه يعلق مساءلة المقاول الأصلي على شروط عسيرة في مقدمتها ثبوت إعسار المقاول من الباطن وعدم تقييده في السجل التجاري وعدم توفره على أصل تجاري، ويحصر الأداء في حدود المبالغ المترتبة بذمته، مما يفرغ الحماية الاجتماعية للأجراء من محتواها العملي.
    \item \textbf{الفرضية الرابعة (رهان الفعالية الرقابية):} إن صيانة التوازن الاجتماعي وحماية العمال في هذه الأنماط المرنة تظل رهينة بتعزيز الصلاحيات التقديرية والضبطية لمفتشي الشغل، وتكريس شجاعة القضاء الاجتماعي في تطبيق مبدأ "أولوية الواقع على المظهر التعاقدي"، فضلاً عن مراجعة العقوبات الزجرية الهزيلة المنصوص عليها في المادة 505 من المدونة.
\end{enumerate}

%% ---------------------------------------------------------------------
\section*{المنهجية المعتمدة وتقسيم التقرير}
%% ---------------------------------------------------------------------

ولمعالجة الإشكالية وفحص الفرضيات، تعتمد هذه الدراسة تكاملاً منهجياً متعدد الروافد:
\begin{itemize}
    \item \textbf{المنهج التحليلي التأصيلي:} لتشريح وتفسير النصوص القانونية المؤطرة في مدونة الشغل (القانون 65.99) والقوانين المكملة (قانون الضمان الاجتماعي وقانون التعويض عن حوادث الشغل)، ورصد مكامن قوتها وثغراتها.
    \item \textbf{المنهج المقارن:} لمقارنة المنظومة القانونية المغربية مع نظيرتها في التشريعات المقارنة (خاصة التشريع الفرنسي) ومعايير منظمة العمل الدولية، للوقوف على التطورات الحديثة في تكييف المسؤوليات.
    \item \textbf{منهج دراسة الحالة والمنهج العملي:} لتنزيل القواعد المجردة على واقع النازلة المعروضة، وتشخيص أدوار الفاعلين الحقيقيين (المشغل، الوسيط، الأجراء، مفتش الشغل، القضاء).
\end{itemize}

وبناءً على هذه المقاربة، جرى تقسيم هذا التقرير البحثي إلى عشرة فصول متناسقة ومبنية وفق هندسة موحدة:
\begin{itemize}
    \item \textbf{الفصل الأول:} الإطار المفاهيمي (الشغل المؤقت، المناولة، التمييز بينهما).
    \item \textbf{الفصل الثاني:} الإطار القانوني في مدونة الشغل المغربية.
    \item \textbf{الفصل الثالث:} الحالات التي يسمح فيها القانون باللجوء إلى الشغل المؤقت.
    \item \textbf{الفصل الرابع:} التزامات شركة التشغيل المؤقت والمقاولة المستعملة.
    \item \textbf{الفصل الخامس:} المسؤولية القانونية لكل من الشركة الأصلية وشركة المناولة.
    \item \textbf{الفصل السادس:} دور مفتش الشغل ومراقبة احترام القانون.
    \item \textbf{الفصل السابع:} الآثار الاجتماعية والاقتصادية (الأجراء، المقاولة، سوق الشغل).
    \item \textbf{الفصل الثامن:} دراسة الحالة وتحليل أدوار الفاعلين.
    \item \textbf{الفصل التاسع:} الاجتهاد القضائي والتجارب المقارنة.
    \item \textbf{الفصل العاشر:} الاختلالات والتوصيات.
\end{itemize}

وينتهي التقرير بخاتمة عامة تركيبية تجيب عن الإشكالية المحورية وتلخص الحلول القانونية لأسئلة النازلة، متبوعة بملحقين (نصوص قانونية مختارة من مدونة الشغل ونماذج عقود استرشادية) ولائحة المراجع المعتمدة.